\section{Experimental Setup}
\label{sec:experimental-setup}

Benchmark, model pool, and protocol are \textbf{not pre-specified} in this draft.
They are selected in \textbf{Phase~A} (research program \S0.7):
\textbf{M1} Experimental Setting Specification,
\textbf{M2} Setting Feasibility Assessment on a selection holdout,
\textbf{M3} Experimental Setting Lock.
Routing opportunity $r(q)$ depends jointly on benchmark, pool, and protocol---these are frozen together in M3.
Table~\ref{tab:setup} is filled once M3 completes.

% T1: Locked experimental configuration
\begin{table}[t]
  \centering
  \small
  \begin{tabular}{ll}
    \toprule
    \textbf{Component} & \textbf{Locked value} \\
    \midrule
    Weak model ($M_w$)   & Llama-3.2-1B-Instruct \\
    Strong model ($M_s$) & Llama-3.2-3B-Instruct \\
    Primary dataset      & ARC-Challenge \citep{clark2018think} \\
    CALIB split          & Official validation ($n{=}299$) \\
    TEST split           & Official test ($n{=}1{,}172$) \\
    Prompt formatting     & Deterministic chat template (\texttt{prompt\_protocol.py}) \\
    Random seed          & 42 \\
    $c(q)$ selection     & \texttt{piece\_count} (CALIB screening) \\
    \bottomrule
  \end{tabular}
  \caption{Locked experimental configuration. Screening exclusions in appendix.}
  \label{tab:setup}
\end{table}


\paragraph{Phase A (M1--M3).}
\textbf{Phase~A exists solely to establish a valid experimental environment. It is not part of the hypothesis-testing pipeline.}
M1 locks candidates, two \emph{deployment scenarios} (\texttt{narrow\_gap}, \texttt{wide\_gap}), protocol layers, evaluation-corpus partition policy, and feasibility gates \emph{before} inference.
M2 runs a factorial pilot over candidate Experimental Settings $\mathcal{S}=(\mathcal{D},\mathcal{P},\Pi)$; selection uses routing buckets only (not signal AUROC).
M3 chooses one primary $\mathcal{S}^*$, assigns split IDs from the locked partition rule, and freezes the Experimental Setting YAML.
Benchmark selection is a reproducibility protocol, not a paper contribution.

\paragraph{Experimental Setting validity.}
An Experimental Setting is \textbf{valid} if it satisfies the requirements necessary to evaluate routing policies under a fixed protocol (Gates A--E).
Validity concerns the experimental environment only and does \textbf{not} imply that the proposed routing signals will be effective.

\paragraph{Evaluation corpus and splits (M1 policy; M3 IDs).}
For each benchmark, M1 declares an \textbf{evaluation corpus} $C$ from the publicly labeled evaluation data (native \texttt{train} excluded; \texttt{dev} used for few-shot only where applicable).
We partition $C$ \textbf{once} into three disjoint subsets with a fixed random seed:
\begin{itemize}
  \item $H$ --- selection holdout (M2 pilot only; never reused);
  \item $R_c$ --- calibration (H1--H3, $(\lambda,\tau)$ fit);
  \item $R_t$ --- test (H4 only).
\end{itemize}
We do \textbf{not} claim official benchmark leaderboard scores; we study routing methodology under a pre-specified partition.
Example (ARC-Challenge): $C =$ validation + test ($|C|=1471$); e.g.\ $|H|=150$, $|R_c|=500$, $|R_t|=821$.

\paragraph{What is tuned.}
Only the routing policy parameters $(\lambda,\tau)$ are fit on $R_c$ (Section~\ref{sec:routing-policy}).
Benchmark, pool, prompt, decoding, grading, and signal definitions are frozen before the main experiment; signal features $x(q)$ are computed, not tuned.

\paragraph{Model pool (locked in M3).}
Fixed pool $\Mpool$ with documented capability ordering; primary pair $\Mlo$, $\Mhi$ with $\mathrm{cap}(\Mlo)<\mathrm{cap}(\Mhi)$; $\kappa = c(\Mhi)/c(\Mlo) > 1$.
Candidate pools are named \textbf{deployment scenarios} (small vs.\ large capability gap), not arbitrary labels.

\paragraph{Dataset and splits (locked in M3).}
Objective correctness $y(q,M)\in\{0,1\}$ under the locked protocol.
\textbf{Calib} $R_c$: signal analysis (H1--H3), selection, $(\lambda,\tau)$ fit (Stages~6--8).
\textbf{Test} $R_t$: routing metrics only (H4, Stage~9).

\paragraph{Oracle labels (Stage 4).}
For each $(q,M)$ in the locked pool we run full inference and record $y(q,M)$, $r(q)$ (from the primary pair), and $\pi^*(q)$ (Section~\ref{sec:problem}).
Oracle labels are used on calib for analysis and policy fit only---never for signal extraction.

\paragraph{Evaluation metrics (test only).}
\begin{align}
  \mathrm{Acc}(\pi) &= \frac{1}{|\mathcal{Q}_{\text{test}}|} \sum_{q\in\mathcal{Q}_{\text{test}}} y\big(q,\pi(q)\big), \\
  \mathrm{Cost}(\pi) &= \frac{1}{|\mathcal{Q}_{\text{test}}|} \sum_{q\in\mathcal{Q}_{\text{test}}} c\big(\pi(q)\big).
\end{align}
We report \textbf{Pareto} curves/tables over $(\mathrm{Acc},\mathrm{Cost})$ for: always-$\Mlo$, always-$\Mhi$, oracle $\pi^*$, best 4a rule, and 4b learned policy.
Optional calib-only tuning uses $J_\alpha(\pi)=\mathrm{Acc}(\pi)-\alpha\,\mathrm{Cost}(\pi)$; reported numbers are always on test.

\paragraph{Baselines.}
Always-$\Mlo$, always-$\Mhi$, oracle $\pi^*$, and our 4a/4b policies from Section~\ref{sec:routing-policy}.
We do not tune $x(q)$ or $(\lambda,\tau)$ on test.
